% Lingua
\usepackage[T1]{fontenc}
\usepackage[italian]{babel}
\usepackage{microtype}
\usepackage[normalem]{ulem}
\usepackage{csquotes}
\usepackage{epigraph}
\setlength{\epigraphwidth}{0.5\textwidth}

% Font
\usepackage{libertinus}
% \usepackage{lmodern}

% Layout
\usepackage{geometry}
\geometry{
  a4paper,
  top=25mm,
  bottom=22mm,
  left=30mm,
  right=30mm
}

\setlength\parindent{0pt} % no spazio quando si va a capo

% Matematica
\usepackage{mathtools}
\usepackage{amsthm}
\usepackage[makeroom]{cancel}
\usepackage{amssymb}
\usepackage{faktor} % per il simbolo del quoziente

% Liste e figures
\usepackage{enumitem}
\usepackage{graphicx}
\usepackage{placeins}

\graphicspath{ {./images/} }

% Bibliografia
\usepackage[
  backend=biber,
  style=alphabetic
]{biblatex}

\addbibresource{references.bib}

% Hyperlinks e numerazione automatica
\usepackage[hidelinks]{hyperref}
\setcounter{tocdepth}{3}

% Colori per i commenti e box colorati
\usepackage{xcolor}
\usepackage[most]{tcolorbox}
\colorlet{myblue}{blue!70!black}
\colorlet{myred}{red!70!black}

\newtcolorbox{mystarthighlight}{
    enhanced, frame hidden,
    colback=blue!3, colupper=myblue,
    left=2pt, right=0pt,
    borderline west={1.5pt}{0pt}{myblue},
}

\newtcolorbox{myendhighlight}{
    enhanced, frame hidden,
    left=8pt, right=0pt,
    borderline west={1pt}{0pt}{black, dashed},
}

% PARAMETRI TCOLORBOX

    % enhanced,                   % abilita le funzionalità grafiche avanzate
    % breakable,                  % permette al box di andare su più pagine

    % COLORI
    % colback=red!3,              % colore dello sfondo
    % colframe=red!75!black,    % colore del bordo normale
    % colupper=black,           % colore del testo interno
    % coltitle=black,           % colore del titolo

    % BORDO
    % frame hidden,               % nasconde il bordo standard
    % borderline={0.8pt}{0pt}{red!75!black,dashed}, % spessore,distanza,colore
    % borderline east = " " "
    % borderline west = " " "
    % borderline nord = " " "
    % borderline sud = " " "
    % boxrule=0.8pt,            % spessore del bordo standard

    % ANGOLI
    % sharp corners,            % angoli completamente quadrati
    % rounded corners,          % angoli arrotondati
    % arc=2mm,                  % raggio degli angoli arrotondati

    % SPAZIATURA INTERNA
    % boxsep=5pt,               % spazio interno generale
    % left=8pt,                 % spazio interno a sinistra
    % right=8pt,                % spazio interno a destra
    % top=6pt,                  % spazio interno sopra
    % bottom=6pt,               % spazio interno sotto

    % SPAZIO ESTERNO
    % before skip=10pt,         % spazio prima del box
    % after skip=10pt,          % spazio dopo il box

    % TITOLO
    % title=Titolo,             % titolo del box
    % fonttitle=\bfseries,      % stile del titolo
    % coltitle=black,           % colore del titolo

    % DIMENSIONI
    % width=\textwidth,         % larghezza del box
    % height=4cm,               % altezza del box

    % ALLINEAMENTO DEL CONTENUTO
    % halign=left,              % allineamento orizzontale (left / center / right)
    % valign=center,            % allineamento verticale (top / center / bottom)

    % EFFETTI GRAFICI
    % drop shadow,              % aggiunge un'ombra al box


% Alcune macro
\newcommand{\Sym}{\text{Sym}}
\newcommand{\gen}[1]{\langle #1 \rangle}
\NewDocumentCommand{\ngen}{m g}{
  \langle\!\langle #1 \rangle\!\rangle_{\IfNoValueTF{#2}{}{#2}}
}
\newcommand{\pres}[2]{\langle #1\mid #2 \rangle}
\newcommand{\mylink}[1]{\href{#1}{[\textbf{link}]}}


%  Numerazione teoremi
\theoremstyle{plain}
\newtheorem{teo}{Teorema}[section]
\newtheorem{prop}[teo]{Proposizione}
\newtheorem{corollario}[teo]{Corollario}
\newtheorem{lemma}[teo]{Lemma}

\theoremstyle{definition}
\newtheorem{defin}[teo]{Definizione}
\newtheorem{esempio}[teo]{Esempio}
\newtheorem{esercizio}[teo]{Esercizio}

\newtheorem{dimostrazione}{Dimostrazione}
\newtheorem{oss}[teo]{Osservazione}


% Personalizzazione documento
